\section{Introduction}

World models compress observations and actions into a state from which future
outcomes can be predicted. Their appeal is clearest when direct interaction is
costly: imagined trajectories can support planning, control, or data-efficient
reinforcement learning. DreamerV3, TD-MPC2, IRIS, and STORM illustrate the
range of recurrent, implicit, and transformer dynamics used in modern agents
\cite{hafner2025_mastering_diverse_control_tasks_through,hansen2023_td_mpc2_scalable_robust_world,micheli2022_transformers_are_sample_efficient_world,zhang2023_storm_efficient_stochastic_transformer_based}.
However, benchmark return does not by itself establish that a transition model
has a reliable forecast origin when recent observations are noisy or missing.
Reviews of model-based reinforcement learning and real-world deployment treat
predictive quality, partial observability, distribution shift, and safety as
distinct problems
\cite{moerland2023_model_based_reinforcement_learning_a,dulac-arnold2021_challenges_of_real_world_reinforcement,brunke2022_safe_learning_in_robotics_from}.

Known dynamics offer useful inductive bias but do not eliminate the forecast-
origin problem. A nominal simulator can be interpretable and data-efficient,
yet parameter drift, unmodeled friction, discretization, and actuator error
make its state estimate biased. A direct measurement is current but may be
corrupted or absent. Physics-guided learning combines mechanistic and neural
components through losses, features, hybrid solvers, and probabilistic state
models
\cite{karniadakis2021_physics_informed_machine_learning,willard2022_integrating_scientific_knowledge_with_machine,yu2024_learning_dynamical_systems_from_data,shen2023_differentiable_modelling_to_unify_machine}.
KalmanNet and data-assimilation research further show that learned or
physics-aware filtering is established prior art
\cite{revach2022_kalmannet,cheng2023_machine_learning_with_data_assimilation}.
The narrower unresolved question is whether a deliberately simple,
prespecified complementary update can provide a more robust starting state for
a long-horizon physics-residual world model under corrupted context.

That question separates three operations that are often evaluated only as a
single end-to-end score. A representation maps the available history into a
state, a context update decides how strongly to trust a new observation, and a
transition propagates the resulting forecast origin. Improving the transition
cannot repair an arbitrarily biased origin, while an accurate context update
cannot prevent a misspecified transition from accumulating error. Missingness
makes the distinction sharper: last-value imputation supplies a numerical
input but does not make that value contemporaneous, whereas nominal propagation
remains current but inherits parameter error. Learned filters can adapt this
trade-off, but their extra flexibility also creates a mechanism-identification
question: an apparent gain over a deterministic physics-guided model may come
from the residual transition rather than from the learned update itself. We
therefore treat equal-trust filtering as a falsifiable engineering choice and
compare it with both no-update and measurement-dominant sensitivities.

We study this question with the Complementary-Filtered Residual World Model
(\method). Within the observed eight-step context, nominal physics predicts the
next state and a fixed gain of 0.5 combines that prediction with each available
measurement; missing channels receive exactly zero measurement gain. The final
filtered state initializes a probabilistic nominal-physics plus bounded neural-
residual transition. This transition is trained by multi-step Gaussian
likelihood and rolled out freely for 25 steps. The fixed update avoids an
unsupported claim that a learned gain is necessary, while gain 0, 0.25, 0.75,
and 1 sensitivities test whether the result is specific to equal trust.

A second gap is evidential rather than architectural. Flexible synthetic
benchmarks make it inexpensive to try many candidate mechanisms, but the same
flexibility can silently turn development outcomes into confirmation. We use a
versioned firewall: two rejected candidates and the normalization defect that
invalidated the first analysis remain archived; a new eight-seed grid was
frozen only after candidate selection; and the runner refuses changes to code,
tuning, methods, seeds, metrics, or decision thresholds. The primary endpoint
tests a practical 5\% margin directly rather than declaring success from an
interval that merely excludes zero. This design makes a negative result as
interpretable as a positive one and prevents favorable system-level cells from
rescuing a failed aggregate hypothesis.

The study is scoped to transition prediction. Five controlled nonlinear
systems---a damped oscillator, Duffing oscillator, Van der Pol oscillator,
torque-driven pendulum, and cart--pole---are evaluated under clean ID,
corrupted ID, and corrupted parameter-shift regimes. The analytical benchmark
makes trajectories, corruption, and shift auditable, but it does not establish
image understanding, policy quality, hardware safety, or real-robot transfer.

Our contributions are:
\begin{enumerate}
 \item a mask-exact complementary context update whose fixed gain and missing-
 channel behavior are frozen before training, coupled to a probabilistic
 physics-residual free rollout;
 \item an exact error decomposition that states when complementary updating can
 reduce forecast-origin error, without claiming universal filter optimality;
 \item an approximately parameter-matched evaluation against five broad
 transition baselines, plus five prespecified gain and training-horizon
 sensitivities over 15 fixed benchmark cells; and
 \item a fail-closed confirmatory design with eight independent seed clusters,
 paired repeated-measurement resampling, direct testing of a 5\% relative-
 improvement margin, attainable multiplicity control, and archived negative
 development decisions.
\end{enumerate}

The primary hypothesis was frozen before any confirmatory seed was generated:
the mean seed-specific NRMSE-25 improvement over a preselected PGNN-style comparator
must be at least 5\%. Support requires both a bootstrap lower bound of 5\% and a
one-sided exact seed sign-flip test against that same margin. This tests the
claimed practical magnitude rather than combining a 5\% point estimate with a
weaker test against zero.

The remainder of the paper positions the method against recent world-model,
physics-guided, and state-estimation work; defines the causal filter and
probabilistic free rollout; specifies the fixed benchmark and inferential
contract; and reports aggregate, regime, comparator, and gain-sensitivity
results. The discussion then separates measured evidence from conditional
mathematical properties and identifies the additional visual, control, and
hardware experiments required for broader embodied-world-model claims.
